\section{Formal Reasoning and Applicability Boundaries}\label{sec:supp-formal}

The method section of the main paper grounds spill-cost-aware liveness shaping in
three formal results: Theorem~\ref{thm:spill-inev} (the spill-inevitability
certificate), Proposition~\ref{prop:belady-margin} (Belady-margin stability), and
Theorem~\ref{thm:area-extra} (the overflow-area approximation). This section
gives their complete proofs and states precisely the conditions under which each
result applies, so that none of them is overclaimed.

The three results delimit the method's scope from complementary angles. The
spill-inevitability certificate diagnoses capacity limitation through a
working-set lower bound taken over the near-optimal set of schedules: when the
certificate holds, the objective shifts from eliminating spills to lowering their
cost. Belady-margin stability explains why the primary optimization degree of
freedom is schedule order rather than eviction policy: under a distance-dominant
condition, different Belady-style eviction rules select the same victim set, so
any difference between them can stem only from order. The overflow-area
approximation theorem then provides a conditional guarantee that the surrogate
objective $\Phi$ tracks spill traffic in the capacity-bound regime. The following
subsections present the complete proofs in turn, keeping ``strictly valid
theorems,'' ``conditional propositions,'' and ``empirical observations'' clearly
separated.

\subsection{Clean/Dirty Cost Asymmetry}\label{sec:supp-clean-proof}

Section~\ref{sec:formulation} of the main paper partitions buffers into clean and
dirty classes and notes that, although the two occupy identical on-chip space,
their spill costs differ by $2\times$. The following observation gives a
quantitative basis for this asymmetric structure; it is the starting point for
every cost-aware decision that follows.

\paragraph{Observation (clean/dirty cost asymmetry).}
Let $x$ and $y$ be two buffers of equal size $s$, with $x$ clean and $y$ dirty.
If a spill replaces $y$ by $x$, the P2 extra traffic of that spill drops from $2s$
to $s$, that is, by exactly one half.

\paragraph{Deriving the cost ratio.}
The asymmetry follows directly from the definition of the evaluation metric: by
\eqref{eq:supp-extra}, the extra traffic of a clean buffer is $s$ and that of a
dirty buffer is $2s$, so their ratio is $1/2$. This is an immediate corollary of
the metric itself and does not depend on any particular scheduler or dataset. In
the synthetic ablation, holding the DAG, the spill count, and the peak fixed, a
clean reserve and a dirty reserve yield extras of $1536$ and $3072$,
respectively, which we use to check that the implementation faithfully reflects
this definition.

\subsection{A Working-Set Certificate over the Near-Optimal Set}\label{sec:supp-working-set-proof}

Fix a topological order $S$ and write $\pos_S(v)$ for the position of node $v$ in
that order. The logical lifetime of a buffer $b$ is
\begin{equation}
  I_S(b)=[\pos_S(a_b),\pos_S(f_b)).
\end{equation}
The logical live working set of cache $c$ at position $t$ is
\begin{equation}
  W_c^S(t)=\sum_{b:\tau(b)=c,\ t\in I_S(b)}s_b,\qquad
  W_c^S=\max_t W_c^S(t).
\end{equation}
For a candidate set of orders $\mathcal S$, let $E^\star(S)$ denote the minimum
P2 extra traffic over all legal address assignments, spill choices, and spill
placements once the underlying node order $S$ is fixed. Given $\epsilon\ge0$,
define the near-optimal set
\begin{equation}
  \mathcal O_\epsilon=\{S\in\mathcal S:\ E^\star(S)\le (1+\epsilon)\min_{S'\in\mathcal S}E^\star(S')\},
\end{equation}
together with the working-set lower bound over that set,
\begin{equation}
  \underline W_c^\epsilon=\min_{S\in\mathcal O_\epsilon} W_c^S.
\end{equation}

Theorem~\ref{thm:spill-inev} of the main paper asserts that if the working-set
lower bound over the near-optimal extra-traffic subset exceeds the cache
capacity, then spilling is unavoidable. We now give the complete proof and
discuss the sense in which this is a one-sided condition: it can confirm that
spilling is unavoidable, but the converse does not hold.

\paragraph{Theorem~\ref{thm:spill-inev} (spill inevitability).}
If $\underline W_c^\epsilon>\Capc_c$, then every $S\in\mathcal O_\epsilon$ must
incur a capacity spill on cache $c$. Equivalently, within the candidate range
that preserves near-optimal P2 extra traffic, address placement alone cannot make
cache $c$ entirely spill-free. More quantitatively, for any $S\in\mathcal
O_\epsilon$ and any legal P2 implementation $P$ respecting $S$,
\begin{equation}
  \extra(P)\ge \max_t (W_c^S(t)-\Capc_c)_+ .
\end{equation}

\paragraph{Proof.}
At any position $t$, decompose the logical live volume into the resident volume
$R_c^P(t)$ and the volume that has been spilled and is momentarily absent,
$A_c^P(t)$:
\begin{equation}
  W_c^S(t)=R_c^P(t)+A_c^P(t).
\end{equation}
The capacity constraint gives $R_c^P(t)\le \Capc_c$, hence
\begin{equation}
  A_c^P(t)\ge (W_c^S(t)-\Capc_c)_+ .
\end{equation}
Any buffer that is momentarily absent yet still live must have been spilled at
least once before this position, and a single spill contributes extra traffic at
least equal to that buffer's size. Therefore the total extra traffic is no
smaller than the absent live volume at any instant. Taking the maximum over $t$
yields the displayed bound. Finally, if $\underline W_c^\epsilon>\Capc_c$, then
every order in $\mathcal O_\epsilon$ has some position at which
$W_c^S(t)>\Capc_c$, so spilling cannot be entirely avoided on that cache.

\paragraph{One-sidedness and scope of the certificate.}
The theorem provides a one-sided certificate: it can confirm only that spilling is
unavoidable, and the converse does not hold. When $\underline W_c^\epsilon\le
\Capc_c$, one cannot conclude spill-freedom; fragmentation, address-reuse
dependencies, pinned resident segments, and reload positions may still force
spills. The restriction to the near-optimal set is also indispensable: if one
minimizes the peak over all legal orders without constraint, certain orders with
very poor P2 extra traffic but low peaks can make an instance appear ``not
capacity-bound.'' For example, the full-order minimum working-set ratio for
Matmul\_Case1/L1 is $0.03125$, whereas within the $10\%$ near-optimal
extra-traffic set the ratio is $8.5$; it is the latter that reflects the capacity
pressure in the P2-relevant region. The CP-SAT oracle validation in
\S\ref{sec:exp-mechanism} of the main paper confirms that this lower bound holds
across all small graphs.

\subsection{A Stability Condition for Belady-Style Eviction}\label{sec:supp-belady-proof}

Proposition~\ref{prop:belady-margin} of the main paper states that, under a
distance-dominant condition, different Belady-style eviction rules select the
same victim set. We now give the formal condition and its proof.

Fix a schedule order. At a cache-overflow decision $k$, let $R_k$ be the set of
evictable resident buffers, $\Delta_k$ the volume that must be freed, and $d_k(b)$
the next-use distance of buffer $b$. Consider a family of distance-dominant
Belady-style rules
\begin{equation}
  \mathrm{score}_r(b)=d_k(b)g_r(b),
\end{equation}
where $g_r(b)>0$ is a bounded, rule-dependent perturbation, such as a clean/dirty
cost weight. Assume that all rules and candidates satisfy
$g_r(b)\in[g_{\min},g_{\max}]$, and write $\rho=g_{\max}/g_{\min}$.

\paragraph{Proposition~\ref{prop:belady-margin} (Belady-margin stability).}
Suppose there is a set $F_k\subseteq R_k$ satisfying: (i) $\sum_{b\in
F_k}s_b\ge\Delta_k$; (ii) no proper subset of $F_k$ suffices to free $\Delta_k$;
and (iii) every $x\in F_k$ and $y\in R_k\setminus F_k$ satisfy
\begin{equation}
  d_k(x)>\rho\,d_k(y).
\end{equation}
Then all the rules above prefer the same victim set $F_k$, possibly differing only
in the tie order within $F_k$. If this condition holds at every overflow
decision, then these rules produce the same residency evolution and the same P2
extra traffic.

\paragraph{Proof.}
For any rule $r$, $x\in F_k$, and $y\notin F_k$,
\begin{equation}
  \frac{\mathrm{score}_r(x)}{\mathrm{score}_r(y)}
  =
  \frac{d_k(x)g_r(x)}{d_k(y)g_r(y)}
  \ge
  \frac{d_k(x)g_{\min}}{d_k(y)g_{\max}}>1.
\end{equation}
Hence every element of $F_k$ ranks ahead of every element outside $F_k$. A greedy
selection that picks the highest-scoring elements until enough volume is freed
therefore selects exactly $F_k$, by sufficiency and minimality.

\paragraph{Intuition and applicability.}
The cost-weighted distance score used in this work is $d_k(b)s_b/e_b$. Since
$e_b=s_b$ or $2s_b$, the clean/dirty weight induces a perturbation of at most
$2\times$. Consequently, when the next-use distance of a far-future candidate
exceeds twice that of a near-term candidate, the cost weight cannot change the
victim set. If the candidates are approximately homogeneous in size and
clean/dirty class, different Belady-style rules likewise degenerate to the same
ordering. This proposition explains why the victim rule is often insensitive on a
fixed order, but it is not an unconditional theorem: if short-lived clean buffers
and long-lived dirty buffers alternate repeatedly, the rule differences can
accumulate. The controlled experiments in \S\ref{sec:exp-mechanism} of the main
paper confirm this prediction: on a fixed order, the extra traffic of four victim
rules differs by $\mathrm{CV}<4\%$, whereas the swing in extra traffic across
different orders can reach $10\times$ or more.

\subsection{Conditionality of the Overflow-Area Surrogate}\label{sec:supp-area-proof}

Let the unnormalized overflow area be
\begin{equation}
  A(S)=\sum_t\sum_c (W_c^S(t)-\Capc_c)_+,
\end{equation}
and the normalized surrogate be
\begin{equation}
  \Phi(S)=\sum_t\sum_c \frac{(W_c^S(t)-\Capc_c)_+}{\Capc_c}.
\end{equation}
Both measure the ``over-capacity memory-time area.'' P2 extra traffic, by
contrast, is an event cost: a single spill can keep a buffer absent across many
schedule positions. Consequently, there is no unconditional constant-factor
approximation between $A(S)$ (or $\Phi(S)$) and the optimal extra traffic. A
long-lived small overflow makes the area grow with time while the extra traffic is
paid only once; conversely, a large buffer that exceeds capacity by only one unit
may nonetheless force the eviction of an entire large block.

Theorem~\ref{thm:area-extra} of the main paper establishes a two-sided constant
relation between the overflow area $A(S)$ and the optimal extra traffic under a
bounded-absence-duration assumption. We now give the complete proof together with
a counterexample construction.

\paragraph{Theorem~\ref{thm:area-extra} (conditional overflow-area approximation).}
Fix an order $S$. Suppose there exists a legal spill plan $P$ whose every spill
episode has absence duration $\ell_j$ satisfying $\ell\le \ell_j\le L$, whose
removed memory-time area
\begin{equation}
  R(P)=\sum_j s_j\ell_j
\end{equation}
satisfies $R(P)\le \lambda A(S)$, and whose every unit spill cost $\kappa_j$
satisfies $\kappa_{\min}\le\kappa_j\le\kappa_{\max}$. Then the fixed-order optimal
extra traffic satisfies
\begin{equation}
  \frac{\kappa_{\min}}{L}A(S)
  \le
  E^\star(S)
  \le
  \frac{\kappa_{\max}\lambda}{\ell}A(S).
\end{equation}
In the evaluation model of this work, $\kappa_{\min}=1$ and $\kappa_{\max}=2$.

\paragraph{Proof sketch.}
Any legal implementation must remove at least $(W_c^S(t)-\Capc_c)_+$ of volume at
each position, so $A(S)\le R(P)$. Since $\ell_j\le L$, we have $R(P)\le L\sum_j
s_j$; as each unit of extra traffic costs at least $\kappa_{\min}$, this yields the
lower bound. For the upper bound we use the feasible plan assumed in the
hypothesis: from $\ell_j\ge\ell$ we obtain $\sum_j s_j\le R(P)/\ell\le \lambda
A(S)/\ell$, which, multiplied by the maximum per-unit cost $\kappa_{\max}$, gives
the upper bound.

\paragraph{From $A$ to $\Phi$.}
Because each cache capacity is fixed, with minimum capacity $256$ and maximum
capacity $4096$,
\begin{equation}
  256\,\Phi(S)\le A(S)\le 4096\,\Phi(S).
\end{equation}
Stating a direct theoretical guarantee in terms of $\Phi$ therefore incurs an
additional $16\times$ capacity-span factor. The main reason this work uses $\Phi$
is not that it is empirically much stronger than per-cache peak pressure, but that
it provides a lifetime-aware, cross-cache-comparable candidate-selection signal at
$O(N)$ cost. The empirical data reported in \S\ref{sec:exp-mechanism} of the main
paper corroborate this positioning: over the evaluation set, the global Spearman
correlation between $\Phi$ and P2 extra traffic is $0.958$, against $0.955$ for
per-cache peak pressure, so the two are essentially on par; and the empirical
ratio of Theorem~\ref{thm:area-extra}, $E^\star/A\in[0.56,1.13]$, falls within the
constant-factor band. This fact also requires restraint in the main paper: the
value of $\Phi$ lies in computational efficiency and cross-stage consistency,
not in predictive accuracy significantly better than per-cache peak.

\subsection{Monotonicity of Portfolio Selection}\label{sec:supp-portfolio-proof}

The method section of the main paper (\S\ref{sec:liveness-shaping}) notes that
adding candidate orders cannot degrade the result. We give the formal argument
here. The safety of the candidate portfolio follows from a simple fact. Let
$\mathcal P$ be a finite set of candidate orders, victim rules, or prefetch
windows, and let $J(x)$ be the true evaluation key (which may be a lexicographic
key). If the output is chosen by $\arg\min_{x\in\mathcal P}J(x)$, then adding a new
candidate $z$ gives
\begin{equation}
  \min_{x\in\mathcal P\cup\{z\}}J(x)\le \min_{x\in\mathcal P}J(x).
\end{equation}
Hence, as long as the final selection uses the true P2/P3 evaluation key rather
than a surrogate score, enlarging the candidate set cannot degrade the selected
result; the surrogate $\Phi$ is used only to generate and screen candidates and
does not replace the final objective.
